\documentclass[12pt]{article}

\usepackage{amsmath}
\usepackage{sbc-template}
\usepackage{natbib}
\usepackage{graphicx,url}
\usepackage{tabularx}
\usepackage[brazil]{babel}
\usepackage[utf8]{inputenc}
\usepackage{algorithm}
\usepackage{algorithmic}
\usepackage{float}
\usepackage{indentfirst}
\usepackage{booktabs}
\usepackage{placeins}
\usepackage[none]{hyphenat}

\emergencystretch=3em

\title{Recuperacao Semantica de Eventos Historicos para Previsao Hibrida de Series Temporais Financeiras}

\author{Jefferson Silva dos Anjos\inst{1}, Luiz Jose Henrique Nogaroli Cavalcante\inst{2}, Eduardo Soares Ogasawara\inst{3}}

\address{\inst{1}Programa de Pos-Graduacao em Ciencia da Computacao (PPCIC), CEFET/RJ\\
\inst{2}Programa de Pos-Graduacao (PPPRO), CEFET/RJ\\
\inst{3}Centro Federal de Educacao Tecnologica Celso Suckow da Fonseca (CEFET/RJ)\\
Rio de Janeiro -- RJ -- Brasil}

\begin{document}

\maketitle

\begin{resumo}
Este artigo propoe um metodo hibrido para previsao financeira de curto prazo que combina modelos quantitativos com recuperacao semantica de eventos historicos comparaveis. O pipeline detecta eventos por limiar de retorno, estrutura noticias em JSON, associa a cada evento sua sequencia realizada de impacto e reutiliza essa biblioteca em avaliacao \textit{walk-forward}. O metodo recupera eventos semanticamente semelhantes anteriores a data prevista, aplica um ajuste por impacto ponderado por similaridade e aprende uma correcao residual sensivel a evento. Experimentos com PETR4, PRIO3 e EXXO34, sincronizados com o Brent, geram 12.159 previsoes fora da amostra sobre baselines LSTM, Autoencoder recorrente e Transformer. O metodo reduz o RMSE medio de 1.6248 para 0.8470 (47,9\%), o MAE medio de 1.2445 para 0.6187 (50,3\%) e supera os modelos base em 76,8\% dos dias avaliados. Os resultados mostram que a recuperacao de eventos comparaveis e um complemento util e auditavel aos modelos puramente quantitativos.
\end{resumo}

\begin{palavras-chave}
Palavras-chaves: Previsao de series temporais financeiras; Recuperacao de eventos; Noticias; Modelos hibridos; Similaridade semantica; Correcao residual.
\end{palavras-chave}

\begin{abstract}
This paper proposes a hybrid method for short-horizon financial forecasting that combines quantitative models with semantic retrieval of comparable historical events. The pipeline detects events through return thresholds, structures news into JSON records, associates each event with its realized post-event sequence, and reuses this event library under walk-forward evaluation. The method retrieves semantically similar past events, applies a similarity-weighted impact adjustment, and learns an event-aware residual correction. Experiments on PETR4, PRIO3 and EXXO34, synchronized with Brent, generate 12,159 out-of-sample forecasts across LSTM, recurrent autoencoder and Transformer baselines. The hybrid method reduces average RMSE from 1.6248 to 0.8470 (47.9\%), reduces average MAE from 1.2445 to 0.6187 (50.3\%), and outperforms the base models on 76.8\% of the evaluated days. These results show that retrieving comparable historical events is a useful and auditable complement to purely quantitative forecasting.
\end{abstract}

\begin{keywords}
Keywords: Financial time series forecasting; Event retrieval; News-aware forecasting; Hybrid models; Semantic similarity; Residual correction.
\end{keywords}

\section{Introducao}

Mercados financeiros reagem continuamente a informacoes corporativas, macroeconomicas e geopoliticas. Embora essa premissa seja bem aceita, a incorporacao operacional de noticias em modelos preditivos ainda permanece incompleta. Em boa parte da literatura, os modelos sao treinados apenas com historico de precos. Quando texto e incorporado, seu uso costuma ocorrer como uma variavel contemporanea de sentimento ou como um canal adicional de classificacao de movimento \cite{Tetlock:2007, Tetlock:2008, Xu:2018, Sawhney:2020}. Esse desenho deixa em aberto uma questao metodologica importante: como reutilizar eventos historicos semanticamente semelhantes e suas trajetorias observadas de impacto para melhorar previsoes numericas de curto prazo?

O problema e particularmente relevante em ativos expostos a choques recorrentes de commodity, regulacao e geopolitica, como os do setor de oleo e gas. Nesses casos, o mesmo tipo de noticia tende a reaparecer em contextos diferentes, com efeitos que nem sempre se esgotam no dia do evento. A abordagem proposta investiga essa lacuna ao converter noticias em eventos estruturados, estimar suas sequencias realizadas de impacto e reaproveitar essas sequencias quando um evento semanticamente comparavel volta a ocorrer.

A abordagem proposta organiza a camada informacional de forma modular. A avaliacao offline principal recupera diretamente pares \textit{motivo-sequencia} da biblioteca historica de eventos, aplica um ajuste por impacto semantico com corte temporal estrito e depois aprende uma correcao residual em regime \textit{walk-forward}. Os clusters semanticos continuam existindo, mas cumprem papel principalmente exploratorio e de auditabilidade.

Nossa hipotese de pesquisa e que recuperar eventos semanticamente comparaveis e suas sequencias observadas de impacto reduz o erro dos modelos base e melhora a auditabilidade da previsao em relacao a abordagens puramente quantitativas.

Com base nisso, as contribuicoes centrais do artigo sao:

\begin{itemize}
    \item uma biblioteca de eventos financeiros estruturados a partir de noticias, com propagacao observada de impacto apos o evento;
    \item um mecanismo de recuperacao semantica de eventos historicos comparaveis com restricao temporal explicita;
    \item um modelo hibrido em duas etapas, composto por ajuste sequencial por similaridade e por correcao residual \textit{walk-forward};
    \item uma avaliacao comparativa sobre tres arquiteturas de base e tres ativos do setor de oleo e gas, com evidencias quantitativas, ablacoes e logs de decisao.
\end{itemize}

\section{Referencial Teorico}

\subsection{Previsao financeira com modelos quantitativos}

Modelos baseados em aprendizado profundo tem sido usados com frequencia na previsao financeira por sua capacidade de capturar relacoes nao lineares e dependencias temporais complexas. Redes LSTM mostraram desempenho competitivo em relacao a abordagens tradicionais \cite{Fischer:2018}, enquanto Autoencoders recorrentes ajudam a construir representacoes latentes mais robustas a ruido \cite{Bao:2017}. Mais recentemente, arquiteturas baseadas em atencao, como o Transformer e suas adaptacoes para series temporais, passaram a disputar esse mesmo espaco por sua eficiencia em sequencias longas e em cenarios multi-horizonte \cite{Vaswani:2017, Lim:2021, Wu:2021}.

\subsection{Noticias, texto e eventos no mercado financeiro}

Ha ampla evidencia de que texto financeiro e fluxo de noticias afetam retorno, volume e percepcao de risco \cite{Tetlock:2007, Tetlock:2008, Loughran:2011, Mitra:2011}. O desenvolvimento de embeddings distribuidos e de modelos contextualizados ampliou a capacidade de representar semelhanca semantica entre frases e eventos \cite{Mikolov:2013, Devlin:2019, Araci:2019, Reimers:2019}. Ainda assim, grande parte dos trabalhos usa o texto como um sinal agregado no instante corrente, e nao como um mecanismo explicito de recuperacao de padroes historicos de impacto.

\subsection{Reacao gradual e memoria de eventos}

Embora a Hipotese dos Mercados Eficientes proponha incorporacao rapida da informacao aos precos \cite{Fama:1970}, a literatura tambem documenta reacoes graduais, exageros e reversoes \cite{Cutler:1989, Barberis:1998, Hong:1999}. Essa linha teorica abre espaco para metodologias que representem o impacto de um evento como uma sequencia temporal observavel, e nao apenas como um choque instantaneo. E precisamente nesse ponto que a recuperacao semantica de eventos historicos se torna relevante: ela oferece uma forma concreta de consultar experiencias anteriores comparaveis e reutilizar suas trajetorias observadas.

\section{Trabalhos Relacionados}

Os trabalhos mais proximos desta proposta podem ser agrupados em quatro familias. A primeira reune modelos puramente quantitativos baseados em series de mercado, como LSTM, Autoencoders e Transformers adaptados para previsao temporal \cite{Fischer:2018, Bao:2017, Lim:2021}. Essas abordagens servem como base de comparacao natural, mas tratam choques informacionais apenas como parte da dinamica endogena da serie.

A segunda familia usa noticias, eventos ou textos de redes sociais para prever direcao de mercado \cite{Ding:2014, Ding:2015, Xu:2018, Sawhney:2020}. O ganho principal desses trabalhos esta em incorporar texto ao processo preditivo, mas em geral o problema e formulado como classificacao de movimento e o texto e fundido ao modelo ponta a ponta, sem gerar uma biblioteca auditavel de sequencias historicas reutilizaveis.

Alem da literatura academica, ha implementacoes publicas amplamente divulgadas que combinam aprendizado profundo, multiplas fontes de atributos e sinais textuais para previsao de mercado. Um exemplo e o repositorio \textit{stockpredictionai}, de Banushev \cite{Banushev:2019}, que organiza uma proposta aplicada com GANs, LSTMs, CNNs e diferentes grupos de \textit{features}. Essa linha e relevante como referencia pratica, mas difere da proposta aqui apresentada por tratar o texto predominantemente como entrada do modelo preditivo, e nao como uma biblioteca explicita de eventos historicos com sequencias observadas e recuperacao auditavel.

A terceira familia avanca para arquiteturas com memoria e interacoes mais estruturadas entre texto, tempo e multiplos ativos \cite{Luo:2023}. Ja a quarta familia aproxima-se de esquemas de recuperacao e memoria baseados em LLMs \cite{Xiao:2025FinSrag, Wang:2025StockMem}. Em ambos os casos, a literatura recente oferece mecanismos sofisticados de memoria, mas tende a depender de arquiteturas mais pesadas, foco em classificacao ou uso dominante do modelo de linguagem como \textit{backbone}.

Em contraste, o presente artigo propoe uma alternativa mais leve e orientada a previsao numerica: a unidade principal de memoria nao e um estado latente do modelo, mas uma biblioteca explicita de eventos historicos com sequencias observadas de impacto; a recuperacao e feita com corte temporal estrito; e a integracao com o preditor ocorre por ajuste sequencial e correcao residual. A lacuna preenchida, portanto, esta na combinacao entre eventos estruturados, recuperacao semantica direta, reutilizacao de sequencias observadas e avaliacao auditavel sobre baselines quantitativos padrao.

\section{Metodologia}

\subsection{Escopo experimental e dados}

O escopo empirico do estudo abrange os ativos PETR4, PRIO3 e EXXO34, sincronizados com o Brent (BZ=F) como proxy de commodity. O pipeline trabalha com aproximadamente cinco anos de dados diarios provenientes do Yahoo Finance, cobrindo o periodo que, na base atual de eventos, vai de 10 de novembro de 2020 a 25 de marco de 2026. O foco setorial foi deliberado: ativos de oleo e gas sao fortemente sensiveis a choques de noticia, commodity, geopolitica e regulacao, o que favorece a observacao de efeitos recorrentes.

\begin{algorithm}[H]
\caption{Resumo do pipeline hibrido proposto}
\label{alg:pipeline}
\small
\begin{algorithmic}[1]
\STATE Detectar dias de evento por limiar de retorno absoluto
\STATE Estruturar cada evento em JSON com contexto textual e motivos
\STATE Construir a biblioteca historica com pares \textit{motivo-sequencia} e datas
\STATE Gerar previsoes base com LSTM, Autoencoder recorrente e Transformer
\FOR{cada data $t$ na avaliacao \textit{walk-forward}}
    \STATE Recuperar eventos anteriores a $t$ com ativo compativel e similaridade semantica acima do limiar
    \STATE Agregar as sequencias recuperadas por media ponderada por similaridade
    \STATE Ajustar a previsao base no horizonte D0--D4
    \STATE Estimar a correcao residual usando apenas a janela passada observada
\ENDFOR
\STATE Registrar previsoes, eventos recuperados e metadados de auditoria
\end{algorithmic}
\end{algorithm}

\subsection{Deteccao e estruturacao de eventos}

Eventos sao disparados quando a variacao percentual diaria absoluta do ativo ultrapassa um limiar predefinido. A versao principal do experimento usa um limiar de 2,0\%, mas a abordagem tambem gera analise de sensibilidade para 1,5\% e 2,5\%. Na configuracao de 2,0\%, a base atual contem 320 eventos para PETR4, 469 para PRIO3 e 301 para EXXO34.

Quando um evento e detectado, o pipeline consulta um modelo da OpenAI para estruturar um registro JSON contendo data, ativo, retorno do dia, descricao do contexto, motivos identificados, sentimento e fontes. Quando o LLM nao encontra contexto suficiente, o pipeline utiliza Tavily como \textit{fallback} de busca e reexecuta a estruturacao do evento sobre o material recuperado. O conjunto experimental armazena atualmente 1.619 arquivos JSON de evento, com 3.253 motivos textuais identificados.

\subsection{Sequencias realizadas e biblioteca de eventos}

Cada evento recebe uma sequencia realizada de impacto apos sua ocorrencia. O processo calcula retornos do dia do evento e dos pregoes seguintes, armazenando sequencias D0--D5 quando houver dados suficientes. Se um novo evento ocorre antes do fim da janela, a propagacao e interrompida, evitando sobreposicao ingenua de impactos. Ao considerar apenas registros com sequencia nao vazia, a biblioteca corrente de recuperacao passa a conter 3.139 pares \textit{motivo-sequencia}.

Formalmente, definimos a biblioteca historica como
\[
\mathcal{L} = \{(m_i, q_i, d_i, a_i)\}_{i=1}^{N},
\]
em que $m_i$ e o motivo textual do evento, $q_i$ e a sequencia observada de impacto, $d_i$ e a data do evento e $a_i$ e o ativo associado. Na inferencia, apenas itens com $d_i < t$ e ativo compativel com a previsao corrente podem ser usados, o que reduz vazamento de informacao futura.

\subsection{Clusters semanticos como camada exploratoria}

Paralelamente a biblioteca direta de eventos, a abordagem constroi uma camada de clusterizacao semantica. Os motivos sao embedados, podem ser ampliados por variacoes textuais quando ha poucos exemplos e sao agrupados por \textit{Agglomerative Clustering}. Na base atual, essa etapa gera 200 clusters exploratorios, 50 para cada um dos quatro grupos (PETR4, PRIO3, EXXO34 e Brent). Esses clusters sao uteis para auditoria qualitativa, para a pergunta de pesquisa sobre recorrencia de padroes e para o MVP online. Na avaliacao offline principal, entretanto, a unidade de recuperacao e o evento historico individual, e nao o cluster agregado.

\subsection{Modelos base e ajuste hibrido}

Tres modelos quantitativos foram treinados como baselines: LSTM, Autoencoder recorrente e Transformer. Todos recebem series sincronizadas entre ativo e Brent e produzem previsoes base de preco. Sobre essa previsao, o metodo hibrido adiciona duas camadas.

No treinamento dos modelos base, os checkpoints sao gerados separadamente para cada combinacao ativo-arquitetura e avaliados diretamente na comparacao offline. O artigo nao realiza selecao \textit{ex post} entre varios checkpoints de uma mesma arquitetura: cada configuracao produz um unico checkpoint final salvo ao termino do treino. As configuracoes principais sao seq\_len igual a 15 para todas as arquiteturas, 140 epocas para LSTM e Autoencoder, 160 epocas para Transformer, taxa de aprendizado de $10^{-3}$ para LSTM e Autoencoder, taxa de $10^{-4}$ para Transformer e, no caso do Autoencoder, peso de reconstrucao igual a 0,5.

A primeira camada e o ajuste por sequencia recuperada. Dado o conjunto de motivos correntes $M_t$, calculamos a similaridade coseno entre seus embeddings e todos os motivos da biblioteca historica filtrada. Mantemos os top-$k$ eventos acima do limiar $\tau = 0{,}70$ e calculamos a sequencia media ponderada
\[
\bar{q}_h = \frac{\sum_{i \in \mathcal{N}_t} s_i q_{i,h}}{\sum_{i \in \mathcal{N}_t} s_i},
\]
em que $s_i$ e a similaridade do item $i$ e $\mathcal{N}_t$ representa o conjunto de eventos recuperados. A previsao base $\hat{y}^{base}_{t+h}$ e ajustada por
\[
\hat{y}^{seq}_{t+h} = \hat{y}^{base}_{t+h}\left(1 + \alpha \cdot \bar{s} \cdot \frac{\bar{q}_h}{100}\right),
\]
em que $\bar{s}$ e a similaridade media dos eventos selecionados e $\alpha$ e um fator de escala escolhido por busca em grade no conjunto $\{0.1, 0.2, 0.3, 0.4, 0.5, 0.6\}$. Na avaliacao offline atual, a recuperacao usa top-$k=5$, limiar de similaridade $\tau = 0{,}70$ e aplica o ajuste sequencial no horizonte D0--D4, ainda que a biblioteca armazene sequencias ate D5.

A segunda camada e uma correcao residual em regime \textit{walk-forward}. Definimos o residual
\[
r_t = y_t - \hat{y}^{seq}_{t},
\]
e treinamos iterativamente um modelo Ridge com atributos derivados apenas do passado observado:
\[
x_t = [r_{t-1}, r_{t-2}, sim_t, event_t, hist_t],
\]
em que $sim_t$ representa a similaridade media do evento corrente, $event_t$ e um indicador binario de evento e $hist_t$ e a quantidade de eventos historicos usados na recuperacao. A previsao final e entao dada por
\[
\hat{y}^{hyb}_{t} = \hat{y}^{seq}_{t} + g(x_t),
\]
com $g(\cdot)$ estimado apenas a partir da janela passada. Na implementacao atual, o ajuste residual usa janela de 15 observacoes. Essa formulacao nao deve ser interpretada como identificacao causal formal; ela representa um mecanismo de ajuste informado por eventos historicos comparaveis sob restricao temporal estrita.

\section{Avaliacao Experimental}

\subsection{Protocolo experimental}

A avaliacao foi realizada de forma offline, em regime \textit{walk-forward}, sobre os tres ativos e as tres arquiteturas base. Para cada combinacao ativo-modelo, o pipeline gera quatro trilhas: previsao base, ajuste apenas por sequencia recuperada (\textit{SeqOnly}), correcao apenas residual (\textit{ResidualOnly}) e previsao hibrida completa. Ao todo, o experimento atual produz 12.159 previsoes fora da amostra, distribuidas uniformemente entre nove configuracoes.

As metricas principais sao RMSE e MAE, por refletirem erro de magnitude em escala original, que e o alvo central deste estudo de regressao. O RMSE recebe maior destaque por penalizar mais fortemente erros grandes em dias de choque, enquanto o MAE complementa a leitura com uma medida mais direta e menos sensivel a outliers. Como metricas complementares, tambem reportamos \textit{directional accuracy}, percentual de dias em que o hibrido supera o baseline e analises de robustez por ano. A \textit{directional accuracy} e mantida como medida auxiliar porque informa acerto de sinal, mas nao captura a magnitude do erro. Alem disso, o pipeline gera logs de decisao por ativo e arquitetura, registrando motivos, similaridade media, quantidade de eventos historicos utilizados e datas de referencia recuperadas. Esses artefatos reforcam a auditabilidade da camada informacional.

Do ponto de vista metodologico, tres cuidados sao centrais no protocolo. Primeiro, a recuperacao historica usa apenas eventos anteriores a data prevista, o que reduz vazamento de informacao futura. Segundo, o ajuste residual e estimado em regime \textit{walk-forward}, com atributos calculados apenas a partir do passado observado. Terceiro, a avaliacao inclui uma ablacao explicita para separar o efeito do ajuste sequencial por similaridade do efeito da correcao residual.

Com esse desenho, a avaliacao permite discutir de forma integrada nao apenas a reducao agregada de erro frente aos baselines quantitativos, mas tambem o comportamento do metodo em dias com e sem evento, a recorrencia de sequencias semanticamente comparaveis, a robustez entre ativos, arquiteturas e recortes anuais, e os limites da abordagem em cenarios de historico pouco comparavel ou sob diferentes limiares de deteccao. Assim, os resultados a seguir sao interpretados menos como respostas isoladas a perguntas independentes e mais como evidencias complementares sobre desempenho, estabilidade, auditabilidade e sensibilidade do metodo proposto.

\subsection{Resultados principais}

A Tabela~\ref{tab:rmse_principal} resume os resultados de RMSE por ativo e arquitetura. O ganho do metodo hibrido aparece em todos os nove cenarios avaliados. A melhora absoluta varia de 0.5167 a 1.4234 pontos de RMSE, com destaque para EXXO34 com Transformer.

\begin{table}[H]
\centering
\small
\caption{RMSE dos modelos base e hibridos por ativo e arquitetura}
\label{tab:rmse_principal}
\begin{tabular}{l l c c c}
\toprule
\textbf{Ativo} & \textbf{Modelo} & \textbf{RMSE Base} & \textbf{RMSE Hibrido} & \textbf{Ganho} \\
\midrule
PETR4  & LSTM         & 0.9632 & 0.4465 & 0.5167 \\
PRIO3  & LSTM         & 1.5592 & 0.8448 & 0.7144 \\
EXXO34 & LSTM         & 2.0610 & 1.1628 & 0.8982 \\
PETR4  & Autoencoder  & 1.0216 & 0.4487 & 0.5730 \\
PRIO3  & Autoencoder  & 1.6077 & 0.8510 & 0.7567 \\
EXXO34 & Autoencoder  & 2.1269 & 1.1725 & 0.9544 \\
PETR4  & Transformer  & 1.0858 & 0.4624 & 0.6234 \\
PRIO3  & Transformer  & 1.3836 & 0.8439 & 0.5397 \\
EXXO34 & Transformer  & 2.8142 & 1.3908 & 1.4234 \\
\bottomrule
\end{tabular}
\end{table}

Em media, o RMSE cai de 1.6248 para 0.8470, o que corresponde a uma reducao agregada de 47,9\%. O MAE medio cai de 1.2445 para 0.6187, uma reducao de 50,3\%. Alem disso, a \textit{directional accuracy} media sobe de 0.4791 para 0.5527, e o metodo hibrido supera o modelo base em 76,8\% dos dias avaliados. Em termos praticos, isso significa que o ganho nao fica restrito a um unico ativo nem a uma unica arquitetura. Ele aparece de forma consistente em LSTM, Autoencoder e Transformer, bem como em empresas mais maduras e mais novas dentro do recorte analisado.

\subsection{Ablacao e interpretacao do ganho}

A Tabela~\ref{tab:ablacao} mostra que a camada \textit{SeqOnly}, isoladamente, gera melhora modesta, enquanto a maior parte do ganho vem da combinacao entre recuperacao semantica e correcao residual \textit{walk-forward}. Isso indica que o valor do metodo nao esta apenas em identificar eventos similares, mas em usar essa informacao como insumo para uma adaptacao residual sensivel ao contexto informacional.

\begin{table}[H]
\centering
\caption{Ablacao agregada por media de RMSE}
\label{tab:ablacao}
\begin{tabular}{l c}
\toprule
\textbf{Configuracao} & \textbf{RMSE medio} \\
\midrule
Modelo base      & 1.6248 \\
SeqOnly          & 1.5789 \\
ResidualOnly     & 0.8881 \\
Hibrido completo & 0.8470 \\
\bottomrule
\end{tabular}
\end{table}

Os numeros da ablacao mostram tres mensagens importantes. Primeiro, ha algum valor na recuperacao semantica pura, mas ela sozinha reduz pouco o erro medio (2,8\%). Segundo, a correcao residual orientada por informacao de evento e o principal motor de ganho, reduzindo o RMSE medio em 45,3\% frente ao baseline. Terceiro, a combinacao das duas etapas ainda acrescenta melhora adicional, levando o ganho agregado a 47,9\%.

\subsection{Resultados por regime e robustez temporal}

Ao separar as observacoes em dias marcados pela biblioteca de eventos e dias sem marcacao, o hibrido melhora em ambos os casos. Em termos medios de MAE, a reducao e de 48,4\% nos dias com evento e de 54,6\% nos dias sem evento. Esse resultado sugere que a camada residual aprende regularidades uteis mesmo fora dos dias explicitamente marcados. Ainda assim, o metodo continua dependente de uma base semantica coerente para capturar os choques mais informacionais.

As analises auxiliares tambem indicam robustez temporal relevante: 60 de 63 subgrupos ano-ativo-modelo melhoram em relacao ao baseline. As tres excecoes aparecem em janelas iniciais de 2020 com apenas tres observacoes. Esse comportamento deve ser interpretado a luz da heterogeneidade da amostra. O recorte inclui empresas mais maduras e series mais longas, como Petrobras, e empresas relativamente mais novas no mercado observado, o que reduz a quantidade de observacoes disponiveis em janelas muito antigas para alguns ativos. Por isso, os poucos casos sem ganho concentram-se em faixas temporais muito curtas e nao indicam degradacao sistematica do metodo.

\subsection{Discussao}

Em conjunto, os resultados sugerem que a biblioteca de eventos recuperados funciona como um mecanismo de memoria externa util para previsao numerica. Seu valor pratico, contudo, depende de duas condicoes: uma base historica minimamente rica e um procedimento de recuperacao sob restricao temporal. Isso ajuda a explicar por que a combinacao entre recuperacao e residual supera tanto o baseline puro quanto o ajuste semantico isolado. O metodo tambem preserva uma propriedade importante para uso academico: cada ajuste pode ser rastreado a eventos historicos concretos, o que facilita analise qualitativa, inspecao de falhas e discussao economica dos resultados.

Como complemento interpretativo, a Tabela~\ref{tab:topicos_impacto} resume os cinco assuntos com maior alteracao absoluta de preco no horizonte futuro entre os clusters com pelo menos cinco eventos. O criterio usado foi a soma absoluta da trajetoria media posterior no horizonte D1--D5, preservando o sinal medio observado. Assim, a tabela nao mede a polaridade semantica intrinseca do assunto, mas sim o comportamento medio dos precos apos eventos daquele tipo.

\begin{table}[H]
\centering
\small
\caption{Assuntos com maior alteracao absoluta de preco no horizonte futuro}
\label{tab:topicos_impacto}
\begin{tabular}{l c c c}
\toprule
\textbf{Assunto representativo} & \textbf{Ativo} & \textbf{$n$} & \textbf{Impacto D1--D5} \\
\midrule
Preocupacoes fiscais e politicas no Brasil sobre o ativo & PETR4 & 9 & +6.2527 \\
Reducao dos riscos geopoliticos no Oriente Medio & BRENT & 8 & -5.3361 \\
Tensoes geopoliticas entre Russia e Ucrania & BRENT & 14 & -5.1211 \\
Sancoes e restricoes ao petroleo russo & BRENT & 12 & -4.3527 \\
Intervencao estatal e politica na gestao da empresa & PETR4 & 7 & -4.0636 \\
\bottomrule
\end{tabular}
\end{table}

\section{Ameacas a Validade}

\subsection{Validade interna}

A principal ameaca a validade interna decorre da dependencia de servicos externos para estruturacao de noticias e geracao de embeddings. Mesmo com prompts controlados, JSON estruturado e \textit{fallback} de busca, ainda existe variacao induzida pelo modelo de linguagem e pela disponibilidade de conteudo recuperado. Alem disso, a abordagem nao realiza identificacao causal formal; portanto, a relacao entre noticia, evento recuperado e ajuste preditivo deve ser interpretada como associacao informacional util, e nao como prova de causalidade economica.

\subsection{Validade de construto}

O construto ``evento relevante'' depende de uma heuristica baseada em limiar de retorno diario. Embora a abordagem inclua analise de sensibilidade para 1,5\%, 2,0\% e 2,5\%, ainda ha risco de que alguns movimentos importantes fiquem fora da base ou de que certos dias entrem como eventos sem corresponder a um choque informacional claro. De forma semelhante, os clusters semanticos devem ser entendidos como apoio exploratorio e nao como validacao definitiva de categorias economicas.

\subsection{Validade externa}

O escopo empirico esta restrito ao setor de oleo e gas e a tres ativos especificos, todos sincronizados com o Brent. Esse desenho foi escolhido por favorecer interpretacao economica e recorrencia de choques informacionais, mas limita a generalizacao imediata para outros setores, mercados ou frequencias temporais. A extensao da abordagem para bancos, varejo, tecnologia ou dados intradiarios ainda precisa ser demonstrada empiricamente.

\subsection{Validade estatistica e de reproducao}

Embora o numero total de previsoes fora da amostra seja expressivo, alguns recortes anuais iniciais possuem poucas observacoes por causa da maturidade diferente entre os ativos analisados. Empresas mais maduras oferecem janelas historicas mais densas, enquanto ativos relativamente mais novos no recorte observado geram menos pontos em anos iniciais. Essa assimetria afeta sobretudo os recortes temporais mais curtos, justamente onde se concentram os poucos casos sem ganho. Alem disso, os artefatos experimentais ja incluem logs, tabelas auxiliares e testes de reproducibilidade, mas ainda nao contemplam uma bateria estatistica mais ampla, como testes formais de significancia, intervalos de confianca por janela ou validacao cruzada temporal em multiplos cortes.

\section{Conclusao}

Este trabalho apresentou um metodo de previsao hibrida para series temporais financeiras no qual noticias sao convertidas em eventos estruturados e reutilizadas por meio de recuperacao semantica de episodios historicos comparaveis. A contribuicao central do artigo esta em combinar uma biblioteca historica de eventos com corte temporal estrito, um ajuste sequencial ponderado por similaridade e uma correcao residual \textit{walk-forward} sobre baselines quantitativos padrao.

As evidencias quantitativas mostram que essa combinacao melhora de forma consistente os tres baselines avaliados. O RMSE medio cai 47,9\%, o MAE medio cai 50,3\%, a \textit{directional accuracy} aumenta e o metodo supera o baseline em mais de tres quartos dos dias avaliados. A ablacao reforca que o ganho nao vem de um unico componente isolado, mas da interacao entre recuperacao semantica e adaptacao residual. Em paralelo, a analise exploratoria dos clusters indica quais assuntos aparecem associados a trajetorias medias posteriores de maior magnitude, preservando uma camada adicional de auditabilidade.

Como proximos passos, destacam-se quatro direcoes. A primeira e reconstruir a base semantica por janelas temporais para reduzir ainda mais o risco de contaminacao historica. A segunda e ampliar a avaliacao para outros setores e mercados. A terceira e incorporar testes estatisticos mais fortes sobre os ganhos observados. A quarta e congelar artefatos nao deterministas para reproducao integral do experimento.

\bibliographystyle{apalike}
\bibliography{sbc-template}

\end{document}
